\chapter{Il solver: costruire, risolvere, interpretare}
\label{cap:pythongurobi}

Questo capitolo è la cassetta degli attrezzi del laboratorio: come si costruisce un
modello con \texttt{gurobipy}, come si fa girare, come si recupera la soluzione e ---
la parte più importante --- come si interpreta l'output. Una versione autonoma di
questa guida è nel file \texttt{GUIDA\_GUROBI.md} nella cartella del corso.

\section{Installazione e licenza}

\begin{lstlisting}[style=output]
python3 -m pip install gurobipy
\end{lstlisting}

Il pacchetto include una licenza dimostrativa (fino a 2000 variabili e 2000 vincoli),
sufficiente per quasi tutti i modelli della dispensa. Per la licenza accademica completa
e gratuita: registrarsi su \url{https://portal.gurobi.com} con l'email istituzionale,
richiedere una \emph{Named-User Academic License} ed eseguire il comando
\texttt{grbgetkey} indicato dal portale (sotto rete di ateneo o VPN).

\section{I cinque passi per costruire un modello}

Riprendiamo il problema di produzione del capitolo~\ref{cap:richiami}:
$\max 30x_A + 50x_B$ con $x_A + 3x_B \le 90$ (ore), $2x_A + x_B \le 80$ (materiale).

\begin{lstlisting}[style=python]
import gurobipy as gp
from gurobipy import GRB

# Passo 1: il contenitore del modello
m = gp.Model("produzione")

# Passo 2: le variabili decisionali (default: continue, >= 0)
xA = m.addVar(name="xA")
xB = m.addVar(name="xB")

# Passo 3: i vincoli (conservare l'oggetto per leggere poi il duale)
v_ore = m.addConstr(xA + 3*xB <= 90, name="ore")
v_mat = m.addConstr(2*xA + xB <= 80, name="materiale")

# Passo 4: la funzione obiettivo
m.setObjective(30*xA + 50*xB, GRB.MAXIMIZE)

# Passo 5 (debug): scrivere il modello in formato leggibile
m.write("produzione.lp")
\end{lstlisting}

\begin{spiegazione}[Le tre regole d'oro di \texttt{gurobipy}]
\begin{enumerate}
  \item \textbf{Variabili indicizzate} con \texttt{addVars}: \\
        \texttt{x = m.addVars(prodotti, mesi, name="x")} crea la famiglia
        \texttt{x[i,t]}; le somme si scrivono \texttt{gp.quicksum(...)} o
        \texttt{x.sum(i, "*")}.
  \item \textbf{Variabili libere} (che possono essere negative, come l'intercetta $b$
        della SVM o la soglia $\eta$ del CVaR): \texttt{m.addVar(lb=-GRB.INFINITY)}.
        Dimenticare \texttt{lb} è l'errore più frequente del laboratorio: il default
        è $\ge 0$!
  \item \textbf{Famiglie di vincoli} con generatore:
        \texttt{m.addConstrs((x.sum("*", t) <= C[t] for t in mesi), name="cap")}.
\end{enumerate}
\end{spiegazione}

\section{Far girare il modello e leggere il log}

\begin{lstlisting}[style=python]
m.Params.OutputFlag = 0    # 0 = silenzioso, 1 = log a video
m.Params.TimeLimit  = 60   # secondi
m.Params.NonConvex  = 2    # serve solo per QP non convessi (es. pricing p*q)
m.optimize()
\end{lstlisting}

Il log del solver, quando attivo, va letto in tre punti:

\begin{lstlisting}[style=output]
Optimize a model with 2 rows, 2 columns and 4 nonzeros  <- dimensioni attese?
Coefficient statistics: Matrix range [1e+00, 3e+00]     <- dati ben scalati?
Optimal objective  1.900000000e+03                      <- esito e valore
\end{lstlisting}

Se i range dei coefficienti spaziano da $10^{-6}$ a $10^{9}$ il modello è mal scalato:
cambiare unità di misura (migliaia di euro, MWh, \dots) prima di cercare altri errori.

\section{Recuperare la soluzione}

\textbf{Controllare sempre lo stato prima di leggere qualunque numero:}

\begin{lstlisting}[style=python]
if m.Status == GRB.OPTIMAL:
    print(m.ObjVal, xA.X, xB.X)
elif m.Status == GRB.INFEASIBLE:
    m.computeIIS(); m.write("conflitto.ilp")   # quali vincoli si contraddicono
elif m.Status == GRB.UNBOUNDED:
    pass  # quasi sempre manca un vincolo o un bound
\end{lstlisting}

\begin{center}\small
\begin{tabular}{llp{7.6cm}}
\toprule
Attributo & Oggetto & Significato \\
\midrule
\texttt{Status} & modello & \texttt{OPTIMAL}=2, \texttt{INFEASIBLE}=3, \texttt{UNBOUNDED}=5, \texttt{TIME\_LIMIT}=9 \\
\texttt{ObjVal} & modello & valore ottimo dell'obiettivo \\
\texttt{X} & variabile & valore della variabile all'ottimo \\
\texttt{RC} & variabile & costo ridotto (solo LP) \\
\texttt{Pi} & vincolo & prezzo ombra / variabile duale (solo LP) \\
\texttt{Slack} & vincolo & scarto: $0$ $\Rightarrow$ vincolo attivo \\
\texttt{SAObjLow/Up} & variabile & intervallo del coefficiente di costo a base invariata \\
\texttt{SARHSLow/Up} & vincolo & intervallo del termine noto in cui vale il prezzo ombra \\
\bottomrule
\end{tabular}
\end{center}

\begin{attenzione}[Il RHS non è sempre quello che avete scritto]
Se un vincolo contiene una costante a sinistra, come
\texttt{quicksum(x[v,t] ...) + base[t] <= C}, Gurobi la sposta nel termine noto: il
RHS memorizzato è \texttt{C - base[t]}. Chi in un'analisi di sensitività scrive
\texttt{constr.RHS = nuovaC} sta in realtà imponendo \texttt{quicksum(x) <= nuovaC},
cioè un'altra cosa. La forma corretta è \texttt{constr.RHS = nuovaC - base[t]}.
Questo errore ci è capitato davvero preparando il capitolo~\ref{cap:ricaricaev}:
il caso di studio lo discute.
\end{attenzione}

\section{Interpretare l'output: la checklist}
\label{sec:checklist}

\begin{enumerate}
  \item Lo stato è \texttt{OPTIMAL}? Se no, fermarsi e diagnosticare
        (\texttt{computeIIS} per l'inammissibilità, file \texttt{.lp} per l'illimitatezza).
  \item Il valore ottimo ha l'ordine di grandezza atteso?
  \item Quali vincoli sono attivi (\texttt{Slack} $=0$)? Sono quelli previsti?
  \item Quanto valgono i prezzi ombra? Quale risorsa conviene potenziare per prima?
        (Ricordare i segni: in un problema di minimo, un vincolo $\le$ ha
        \texttt{Pi} $\le 0$.)
  \item La soluzione è stabile? Rieseguire con dati perturbati del 5\%.
  \item Tradurre tutto in una raccomandazione manageriale di tre righe.
\end{enumerate}

\begin{esempio}[Lettura completa dell'esempio $2\times2$]
Eseguendo il codice della sezione precedente:
\begin{lstlisting}[style=output]
Status: 2 (OPTIMAL)      ObjVal: 1900.0
xA = 30.0  xB = 20.0
ore:       Slack = 0.0   Pi = 14.0   SARHS = [40.0, 240.0]
materiale: Slack = 0.0   Pi = 8.0    SARHS = [30.0, 180.0]
\end{lstlisting}
Lettura manageriale: entrambe le risorse sono sature; un'ora di straordinario vale fino
a 14~\euro{} e la valutazione resta valida finché le ore disponibili restano tra 40 e
240; oltre, va ricalcolato tutto. Se si potesse potenziare una sola risorsa, meglio
le ore (14 $>$ 8 per unità, a parità di costo di espansione).
\end{esempio}

\section{E per i modelli non lineari generali: scipy}

Gurobi copre LP, QP (anche non convessi) e vincoli quadratici. Per gli NLP generali
della dispensa (localizzazione di Weber, domanda logistica, code M/M/1) usiamo la
funzione \texttt{minimize} di \texttt{scipy.optimize}: stessa logica --- dati, funzione obiettivo,
risoluzione, lettura --- ma il risultato è un ottimo \emph{locale}, salvo convessità
dimostrata a parte.

\begin{lstlisting}[style=python]
from scipy.optimize import minimize
res = minimize(f, x0=punto_iniziale, method="SLSQP",
               bounds=[(0, None)] * n,
               constraints=[{"type": "ineq", "fun": g}])  # g(x) >= 0
assert res.success
print(res.x, res.fun)
\end{lstlisting}

\begin{esercizio}
Implementare l'esempio $2\times2$, stampare \texttt{Pi}, \texttt{Slack},
\texttt{SARHSLow/Up} di entrambi i vincoli e verificare il prezzo ombra delle ore
ri-ottimizzando con $90 \to 91$.
\end{esercizio}

\begin{esercizio}
Togliere il vincolo del materiale e osservare l'esito. Poi cambiare l'obiettivo in
minimizzazione senza altri cambiamenti: perché il risultato è $x = 0$?
\end{esercizio}

\begin{esercizio}
Costruire apposta un modello inammissibile (domanda superiore alla capacità totale in
un problema di trasporto) e usare \texttt{computeIIS} per individuare il conflitto.
\end{esercizio}
